\subsection{MMU 相关异常处理}

\subsubsection{异常类型}

MMU相关的异常包括：

\begin{table}[H]
\centering
\caption{MMU相关异常类型}
\begin{tabular}{|c|c|c|c|}
\hline
\textbf{异常名} & \textbf{Ecode} & \textbf{触发条件} & \textbf{处理程序} \\
\hline
TLB重填异常(TLBR) & 0x3F & TLB未命中 & 重填异常处理程序 \\
\hline
Load页无效异常 & 0x08 & Load时V=0 & 普通异常 \\
\hline
Store页无效异常 & 0x09 & Store时V=0 & 普通异常 \\
\hline
取指页无效异常(PIF) & 0x0C & 取指时V=0 & 普通异常 \\
\hline
页特权不合规异常(PPI) & 0x0B & 权限检查失败 & 普通异常 \\
\hline
页修改异常(PME) & 0x0A & Store到D=0页 & 普通异常 \\
\hline
\end{tabular}
\end{table}

\textbf{异常检测逻辑的实现}

在EX级和IF级分别进行异常检测。EX级处理load/store异常，IF级处理取指异常。

\begin{lstlisting}[language=Verilog, caption=EX级地址转换异常检测]
// In EX stage for load/store operations
if (tlb_lookup_result.found == 1'b0) begin
    // TLB miss
    ex_ecode = `ECODE_TLBR;
    ex_excep = 1'b1;
end else begin
    // TLB hit, check various conditions
    case (1'b1)
        !tlb_lookup_result.v: begin
            // Invalid page
            if (is_load) begin
                ex_ecode = `ECODE_PIL;  // Load Invalid
            end else if (is_store) begin
                ex_ecode = `ECODE_PIS;  // Store Invalid
            end
            ex_excep = 1'b1;
        end
        
        (tlb_lookup_result.plv < crmd_plv): begin
            // Privilege violation
            ex_ecode = `ECODE_PPI;
            ex_excep = 1'b1;
        end
        
        (is_store && !tlb_lookup_result.d): begin
            // Page modification exception
            ex_ecode = `ECODE_PME;
            ex_excep = 1'b1;
        end
        
        default: begin
            // No exception
            ex_excep = 1'b0;
        end
    endcase
end
\end{lstlisting}

\begin{lstlisting}[language=Verilog, caption=IF级取指异常检测]
// In Pre-IF stage for instruction fetch
if (tlb_lookup_result.found == 1'b0) begin
    // TLB miss
    if_ecode = `ECODE_TLBR;
    if_excep = 1'b1;
end else if (!tlb_lookup_result.v) begin
    // Invalid page for fetch
    if_ecode = `ECODE_PIF;
    if_excep = 1'b1;
end else if (tlb_lookup_result.plv < crmd_plv) begin
    // Privilege violation
    if_ecode = `ECODE_PPI;
    if_excep = 1'b1;
end else begin
    if_excep = 1'b0;
end
\end{lstlisting}

\textbf{BADV和TLBEHI的自动更新}

发生地址相关异常时，硬件自动将坏虚地址记录到BADV寄存器，并将虚地址的高位记录到TLBEHI的VPPN域。

\begin{lstlisting}[language=Verilog, caption=BADV和TLBEHI的异常更新]
// In WB stage, detect address-related exceptions
wire badv_ex = wb_ex && (
    wb_ecode == `ECODE_ADEF ||  // Address error on fetch
    wb_ecode == `ECODE_ALE  ||  // Address error on load/store
    wb_ecode == `ECODE_PIL  ||  // Page invalid on load
    wb_ecode == `ECODE_PIS  ||  // Page invalid on store
    wb_ecode == `ECODE_PIF  ||  // Page invalid on fetch
    wb_ecode == `ECODE_PME  ||  // Page modification
    wb_ecode == `ECODE_PPI  ||  // Privilege not sufficient
    wb_ecode == `ECODE_TLBR     // TLB refill
);

always @(posedge clk) begin
    if (badv_ex) begin
        // Update BADV with bad virtual address
        csr_badv <= badv_wdata;
        
        // Update TLBEHI VPPN with virtual page number
        csr_tlbehi[`CSR_TLBEHI_VPPN] <= badv_wdata[31:13];
    end
end
\end{lstlisting}

关键点：

\begin{enumerate}
    \item \textbf{准确的坏地址捕获}：对于load/store指令，坏地址应该是计算出的访存虚地址；对于取指，应该是取指请求的虚地址。

    \item \textbf{VPPN自动提取}：硬件从坏地址的[31:13]位自动提取虚拟页号，无需软件干预。
\end{enumerate}

\textbf{TLB重填异常的特殊处理}

TLB重填异常有特殊的处理流程：

\begin{enumerate}
    \item \textbf{自动进入直接翻译模式}：触发TLB重填异常时，硬件自动将CSR.CRMD的DA置1、PG置0，进入直接地址翻译模式。这样异常处理程序自身的虚实转换就不会再触发TLB异常。
    
    \item \textbf{独立的异常入口}：TLB重填异常的入口地址由CSR.TLBRENTRY指定，不是CSR.EENTRY。
    
    \item \textbf{硬件TLB填充}：异常处理程序通过软件页表遍历找到对应页表项，写入TLBEHI、TLBELO0、TLBELO1和TLBIDX寄存器，然后执行TLBFILL或TLBWR指令将其写入TLB。
\end{enumerate}

\begin{lstlisting}[language=Verilog, caption=TLB重填异常处理]
always @(posedge clk) begin
    if (rst) begin
        csr_crmd[`CSR_CRMD_DA] <= 1'b1;
        csr_crmd[`CSR_CRMD_PG] <= 1'b0;
    end
    else if (wb_ex && wb_ecode == `ECODE_TLBR) begin
        // Entering TLB refill exception
        // Automatically switch to direct translation
        csr_crmd[`CSR_CRMD_DA] <= 1'b1;
        csr_crmd[`CSR_CRMD_PG] <= 1'b0;
        
        // Set exception state for TLBFILL/TLBWR to recognize
        csr_estat[`CSR_ESTAT_ECODE] <= `ECODE_TLBR;
    end
    else if (ertn_inst && csr_estat[`CSR_ESTAT_ECODE] == `ECODE_TLBR) begin
        // Exiting TLB refill exception
        // Restore mapped translation mode
        csr_crmd[`CSR_CRMD_DA] <= 1'b0;
        csr_crmd[`CSR_CRMD_PG] <= 1'b1;
    end
end
\end{lstlisting}

% 流水线集成与虚实地址转换流程见integration.tex

\subsubsection{重取指机制的完整流程}

对于影响地址翻译的CSR修改或TLB修改，采用重取指机制：

\begin{enumerate}
    \item \textbf{设置重取指标志}：当WB级检测到需要重取指的指令时，设置全局重取指标志。
    
    \item \textbf{标记流水线指令}：ID级发现重取指标志置位时，将新取来的指令标记为"需要重取"。
    
    \item \textbf{处理重取指指令}：带有重取指标记的指令如同异常一样处理，在WB级时清空整个流水线。
    
    \item \textbf{重新取指}：清空流水线后，PC保持在重取指指令处，从该指令重新开始取指。
\end{enumerate}

% TODO: fill it with implementation details

% \begin{lstlisting}[language=Verilog, caption=重取指机制的完整流程]
% // WB stage: detect instructions needing refetch
% wire wb_refetch_trigger = (wb_is_tlbwr | wb_is_tlbfill | wb_is_invtlb) |
%                           (wb_csr_we & (
%                               (csr_num == `CSR_CRMD) |
%                               (csr_num == `CSR_DMWO) |
%                               (csr_num == `CSR_DMW1) |
%                               (csr_num == `CSR_ASID)
%                           ));
% 
% // Set refetch flag
% always @(posedge clk) begin
%     if (rst || flush) begin
%         refetch_flag <= 1'b0;
%     end
%     else if (wb_refetch_trigger) begin
%         refetch_flag <= 1'b1;
%     end
%     else if (id_refetch_mark_set) begin
%         // Keep flag until refetch instruction reaches WB
%         refetch_flag <= refetch_flag;
%     end
% end
% 
% // ID stage: mark instruction with refetch flag
% wire id_inst_refetch = refetch_flag;
% 
% // WB stage: process refetch instruction
% if (id_inst_refetch) begin
%     // Clear pipeline
%     flush <= 1'b1;
%     
%     // Keep PC at refetch instruction for re-fetch
%     nextpc = current_pc;
%     
%     // Clear refetch flag
%     refetch_flag <= 1'b0;
% end
% \end{lstlisting}

% \subsection{小结}
% 
% TLB虚实地址转换与异常支持的关键实现要点：
% 
% \begin{itemize}
%     \item \textbf{三级翻译流程}：直接翻译 → DMW映射 → TLB映射，按优先级进行
%     
%     \item \textbf{并行异常检测}：在Pre-IF和EX级同时进行异常检测，及早发现问题
%     
%     \item \textbf{完善的CSR管理}：BADV、TLBEHI等CSR自动更新坏地址信息
%     
%     \item \textbf{特殊异常处理}：TLB重填异常自动进入直接翻译模式，有独立入口
%     
%     \item \textbf{重取指机制}：解决配置变更对流水线的影响
% \end{itemize}
